\section{Correctness}
\label{sec:correctness}

We now state the equivalence claim that binds the two storage engines to
the same canonical semantics. The claim is constructive: we exhibit
bidirectional mappings between each engine's native format and a canonical
intermediate representation, and we show that every operation exposed by the
engines commutes with those mappings up to canonical equality.

\subsection{Canonical Fact and Canonical Result}
\label{sec:correctness:canonical}

Let $\mathfrak{F}$ denote the type of \emph{canonical facts}: a canonical fact
is a semantic fact whose fields have been normalised (identifiers sorted,
object tuples canonicalised, metadata keys ordered) so that two facts are
canonically equal exactly when they carry the same semantic content.
Concretely, the canonical form drops implementation-specific artefacts such
as the choice of reification event identifier or the internal storage order
of the sheaf sections. For a query $Q$ and a storage engine $\Sigma$, the
canonical result $\mathfrak{R}_\Sigma(Q)$ is the set of canonical facts
obtained by canonicalising each fact in $\llbracket Q \rrbracket_\Sigma$.

\subsection{Lossless Mapping to Each Engine}
\label{sec:correctness:lossless}

\begin{definition}[Lossless Mapping]
\label{def:lossless}
A mapping $\phi: \mathcal{F} \to \mathcal{S}$ into an engine-specific
representation $\mathcal{S}$ is \emph{lossless} if there exists a reverse
mapping $\psi: \mathcal{S} \to \mathcal{F}$ with $\psi \circ \phi = \mathrm{id}_\mathcal{F}$.
\end{definition}

The implementation provides two such maps:
\begin{itemize}
    \item $\phi_{\mathrm{KG}}: \mathcal{F} \to \mathcal{T}$ reifies a fact
          into a set of typed triples over an auxiliary event node
          (Section~\ref{sec:mapping:kg-real});
    \item $\phi_{\mathrm{Sh}}: \mathcal{F} \to \mathcal{S}_{\mathrm{Sh}}$
          stores a fact as a local section attached to its context
          (Section~\ref{sec:mapping:sh-real}).
\end{itemize}
Both maps have inverses $\psi_{\mathrm{KG}}$ and $\psi_{\mathrm{Sh}}$ that
reconstruct the original fact bit-for-bit; the round-trip is checked by the
test suite for every fact inserted during a benchmark run.

\subsection{Semantic Equivalence of Query Results}
\label{sec:correctness:equivalence}

\begin{definition}[Semantic Equivalence]
\label{def:semanticequiv}
Two storage engines $\Sigma_1, \Sigma_2$ built from the same finite fact set
$F$ are \emph{semantically equivalent} if for every query $Q \in \mathcal{Q}$
we have $\mathfrak{R}_{\Sigma_1}(Q) = \mathfrak{R}_{\Sigma_2}(Q)$.
\end{definition}

\begin{theorem}[Cross-Engine Equivalence]
\label{thm:equivalence}
The KG and SFDB engines built from the same finite $F \subseteq \mathcal{F}$
are semantically equivalent.
\end{theorem}

\begin{proof}[Proof sketch]
Let $\pi: \mathcal{F} \to \mathfrak{F}$ denote canonicalisation. It suffices
to show that $\pi \circ \psi_{\mathrm{KG}} \circ \phi_{\mathrm{KG}} = \pi$ and
$\pi \circ \psi_{\mathrm{Sh}} \circ \phi_{\mathrm{Sh}} = \pi$ on
$\mathcal{F}$, and that both engines evaluate every query $Q$ by scanning
their respective representations of $F$ and returning
$\{\psi(x) \mid x \in \Sigma, x \models Q\}$. The first two identities follow
from losslessness (Definition~\ref{def:lossless}), and the second is the
correctness property of the two query evaluators. Composing the two identities
with $\pi$ on either side gives
$\mathfrak{R}_{\Sigma_{\mathrm{KG}}}(Q) = \mathfrak{R}_{\Sigma_{\mathrm{Sh}}}(Q)$
for every $Q$.
\end{proof}

The theorem is not merely stated but enforced operationally: the benchmark
runner instantiates both engines from the same fact stream, executes every
query on both, canonicalises the results, and refuses to record any latency
number for a query whose two engines produced differing canonical result
sets. Every benchmark row reported in Section~\ref{sec:evaluation} has passed
this check.

\subsection{Referential and Consistency Invariants}
\label{sec:correctness:integrity}

\begin{definition}[Referential Integrity]
\label{def:refintegrity}
A storage engine $\Sigma$ satisfies \emph{referential integrity} if, for
every fact $f = (\mathrm{id}, s, r, \vec{o}, c, k, m) \in \Sigma$, the subject
$s$ is a known entity in $\Sigma$, every non-literal object $\omega_j \in \vec{o}$
is a known entity in $\Sigma$, and the relation $r$ is a known relation in
$\Sigma$.
\end{definition}

Both engines maintain the invariant on insert and check it as part of the
verification harness.

\begin{definition}[Consistency]
\label{def:consistency}
A storage engine is \emph{consistent} if every pair of overlapping sections
agrees under restriction to their common sub-context. Formally, for every
$c, d \in \mathcal{C}$ and every $f \in F(c)$, $g \in F(d)$ such that
$\pi(f) = \pi(g)$ on the intersection $c \wedge d$, we have
$\rho_{c, c \wedge d}(f) = \rho_{d, c \wedge d}(g)$.
\end{definition}

As noted in Section~\ref{sec:sheaf-remark}, on the finite Alexandrov topology
this condition is vacuous: the intersection $c \wedge d$ always exists in the
poset, restriction maps are total functions, and the check reduces to a
tautology. The consistency checker in the implementation therefore serves as
a defensive assertion rather than as a non-trivial theorem-proving procedure.
